\documentclass[12pt, a4paper]{article}
\setlength{\headheight}{14.5pt}
% --- Essential Packages ---
\usepackage[margin=1in]{geometry}
\usepackage{amsmath, amssymb, amsthm}
\usepackage{graphicx}
\usepackage{setspace}
\usepackage{float}
\usepackage{lastpage}

% --- Use Times font (compatible with pdfLaTeX on Overleaf) ---
\usepackage{times}
\usepackage{mathptmx}

% --- Colors and Tables ---
\usepackage[table, xcdraw, dvipsnames]{xcolor}
\usepackage{multirow}
\usepackage[framemethod=TikZ]{mdframed}
\usepackage{enumerate}
\usepackage[shortlabels]{enumitem}

% --- Algorithms and Figures ---
\usepackage[linesnumbered, ruled, vlined]{algorithm2e}
\usepackage{subcaption}

% --- Headers and Footers ---
\usepackage{fancyhdr}
\pagestyle{fancy}
\fancyhf{}
\fancyhead[L]{SWE 4304: Software Project Lab II}
\fancyhead[R]{Project Proposal (Group 6)}
\renewcommand{\footrulewidth}{0.4pt}
\fancyfoot[C]{Page \thepage\ of \pageref{LastPage}}

% --- Hyperlinks (loaded last) ---
\usepackage[colorlinks]{hyperref}
\hypersetup{
    colorlinks=true,
    linkcolor=black,
    filecolor=black,
    urlcolor=black,
    urlbordercolor=blue
}

\begin{document}

\begin{titlepage}
    \begin{figure}[!h]
    \centering
    
\includegraphics[width=.2\textwidth]{iut_logo.png}
    \end{figure}
\begin{center}

\noindent\rule{\textwidth}{1.5pt}
\Huge{\textbf{Islamic University of Technology}}\\
\large{\textbf{Department of Computer Science and Engineering (CSE)}}\\
\large{{B.Sc.  in Software Engineering}}\\
\vspace{0.2cm}
\Large{\textbf{SWE 4304:} Software Project Lab II}\\
\Large{{Summer}, 2023-2024}\\
\noindent\rule{\textwidth}{0.5pt}

\vspace{1cm}
\Large{\textbf{Project Proposal}} \\
\Large{{Plastic\_box\_solver}} \\
\vspace{1cm}

\vspace{0.5cm}
\large{\textbf{Group 6}} \\
\large{S. M. Tahsinuzzaman Emon, 230042104}\\
\large{S. M. Samiul Hossain, 230042110}\\
\large{Md. Muntahi Hasan Akhiar, 230042118}\\
\vspace{1.5cm}
\large{\textbf{Supervisor}} \\
\large{Aashnan Rahman\\ Junior Lecturer, CSE}\\

\vspace{1.5cm}
\large{{Date: \today}} \\

\end{center}
\end{titlepage}

\newpage

\tableofcontents

\newpage

\begin{center}
    \Large{\textbf{Plastic\_box\_solver}}
\end{center}

\section{Project Overview}

Plastic\_box\_solver is a web-based Rubik's Cube assistant that lets users scan their scrambled cube using a phone camera, then either solve it themselves on an interactive 3D model or follow an animated step-by-step guide powered by the Kociemba algorithm. The app runs entirely in the browser, tracks solve time and efficiency, and maintains a global leaderboard so users can compare their progress.

\vspace{0.5cm}
\section{Motivation Behind the Project}

Most people who own a Rubik's Cube never learn to solve it because available resources — video tutorials or raw algorithm sheets — are difficult to follow alongside a physical cube. A browser-based tool that scans the cube, renders it in 3D, and walks the user through each move interactively removes that barrier entirely. Beyond learning, competitive solvers have a genuine need to measure their efficiency against the optimal solution; having a live timer, move counter, and leaderboard turns a solitary puzzle into an engaging, trackable activity.


\vspace{0.5cm}
\section{User Requirement Analysis}

\subsection*{Stakeholder}

One of our friends, who recently got into Rubik's Cube solving, mentioned that he constantly struggles with following YouTube tutorials while also handling his physical cube. He expressed that he would find it far more convenient if an app could simply scan his cube and show him, on an animated model, exactly which moves to make next — without having to switch back and forth between a video and the physical cube.

\subsection*{Functional Requirements}

\begin{itemize}[noitemsep, topsep=5pt]
    \item The system shall allow users to register, log in, and maintain a personal profile.
    \item The system shall accept a scrambled cube state via camera scan or manual color input.
    \item The system shall detect and validate all 54 sticker colors from the submitted images.
    \item The system shall render the cube as a fully interactive, rotatable 3D model in the browser.
    \item The system shall compute an optimal solution using the Kociemba two-phase algorithm.
    \item The system shall support an \textit{Interactive Solve} mode where users manipulate the cube themselves.
    \item The system shall support a \textit{Guided Solve} mode with animated move-by-move instructions.
    \item The system shall track solve time, move count, and efficiency percentage per session.
    \item The system shall persist user history and display a global leaderboard.
\end{itemize}

\subsection*{Non-Functional Requirements}

\begin{itemize}[noitemsep, topsep=5pt]
    \item Color detection shall be accurate under standard indoor lighting.
    \item Solution computation shall return a result within two seconds.
    \item The application shall be responsive and usable on mobile browsers without installation.
\end{itemize}

\newpage
\section{Flow Chart / Class Diagram}

\begin{figure}[H]
    \centering
    \includegraphics[width=0.95\textwidth]{flowchart.png}
    \caption{Application Flow Chart}
    \label{fig:flowchart}
\end{figure}

\newpage
\section{Key Features}

\subsection*{Cube Input}
The user scans each face of their cube using a phone camera or webcam. The Python backend processes each image with OpenCV to detect the nine sticker colors per face and returns a 2D preview for confirmation. If lighting conditions make detection unreliable, a manual click-input interface lets the user set each sticker color directly on a flat 2D cube net.

\subsection*{Interactive 3D Cube}
The cube is rendered in the browser using React Three Fiber, built from 26 individual cubie meshes. OrbitControls let the user drag to rotate the view freely. Face rotations triggered by on-screen buttons (U, R, F, etc.) are smoothly animated, and the 3D model stays in sync with the logical cube state at all times.

\subsection*{Solve Modes}
In \textbf{Interactive Mode}, the user solves the cube themselves; a timer and move counter run live, and the optimal move count is shown on completion for comparison. In \textbf{Guided Mode}, the Kociemba-generated move list appears in a sidebar and the user can step through moves one at a time or trigger Auto-Solve to watch the full solution animate automatically.

\subsection*{Statistics and Leaderboard}
The sidebar shows current solve time, moves taken, and efficiency percentage in real time. All completed solves are saved to the database. Users can view their personal best and compare against others on a global leaderboard ranked by solve time and efficiency.

\vspace{0.5cm}
\section{Tools and Technologies}

\vspace{0.3cm}
\subsection*{Technology Stack}
\begin{table}[H]
    \centering
    \caption{List of Tools and Technologies}
    \label{tab:tools}
    \begin{tabular}{|l|l|}
        \hline
        \rowcolor[HTML]{EFEFEF}
        \textbf{Category} & \textbf{Tool / Technology} \\
        \hline
        Frontend Framework      & React + Vite \\
        \hline
        3D Rendering            & React Three Fiber, @react-three/drei \\
        \hline
        Styling                 & Tailwind CSS \\
        \hline
        State Management        & Zustand \\
        \hline
        HTTP Client             & Axios \\
        \hline
        Backend Framework       & FastAPI (Python) \\
        \hline
        Computer Vision         & OpenCV + NumPy \\
        \hline
        Solving Algorithm       & kociemba \\
        \hline
        ORM                     & SQLAlchemy \\
        \hline
        Database                & PostgreSQL \\
        \hline
        API Testing             & Postman \\
        \hline
        Version Control         & Git \& GitHub \\
        \hline
    \end{tabular}
\end{table}

\vspace{0.5cm}
\section{Proposed Timeline}

\begin{figure}[H]
    \centering
    \includegraphics[width=1\linewidth]{timeline.png}
    \caption{Timeline Diagram (Gantt Chart)}
    \label{fig:timeline}
\end{figure}

\newpage
\section{Action Taken Based on the Feedback}

The suggestions that we received are:
\begin{itemize}[noitemsep, topsep=5pt]
    \item \textbf{Suggestion 1:}
    \item \textbf{Suggestion 2:}
    \item \textbf{Suggestion 3:}
\end{itemize}

\vspace{0.5cm}
\section{Links}
\begin{enumerate}[noitemsep, topsep=5pt]
    \item \href{https://www.canva.com}{Presentation Slide Link}
    \item \href{https://github.com/SMTEmon/plastic-box-solver.git}{GitHub Repository Link}
\end{enumerate}

\end{document}